\documentclass[12pt]{article}
\usepackage[utf8]{inputenc}
\usepackage[T1]{fontenc}
\usepackage[margin=1in]{geometry}
\usepackage{setspace}
\usepackage{xcolor}
\usepackage[colorlinks=true, urlcolor=blue!60!black, linkcolor=blue!60!black]{hyperref}
\setlength{\parindent}{2em}
\setlength{\parskip}{1.25em}
\renewcommand{\baselinestretch}{1.0}
\begin{document}
\begin{center}
{\LARGE \textit{Sherbert v.\ Verner}, 374 U.S.\ 398 (1963)\par}
\vspace{0.8em}
{\large Case Brief --- Four-Page Edition\par}
\end{center}
\vspace{1em}

\section{Memory Jogger}

South Carolina violated the Free Exercise Clause by denying unemployment benefits to a Seventh-day Adventist who would not work on her Saturday Sabbath; substantial burdens on religion require a ``compelling state interest'' (\href{https://supreme.justia.com/cases/federal/us/374/398/}{Justia}).

\section{Detailed Case Facts}

Adell Sherbert worked about 35 years as a spool tender at Spartan Mills in Spartanburg, South Carolina. In 1957 she joined the Seventh-day Adventist Church, whose tenets forbid labor from sundown Friday to sundown Saturday; at the time, Saturday work was optional. In 1959 the mill moved to a mandatory six-day week, and after she missed six successive Saturdays it discharged her. She sought work at three other mills and in other industries but could find no five-day job. The Unemployment Compensation Act required a claimant to be ``able to work and is available for work'' and disqualified one who failed ``without good cause'' to accept ``available suitable work.'' A separate statute protected Sunday worshippers in textile plants from Sunday work during national emergencies.

\section{Procedural History}

The Employment Security Commission, chaired by Charlie V.\ Verner, held that her refusal of Saturday work was a failure without good cause to accept suitable work. The Spartanburg County Court of Common Pleas affirmed, and the South Carolina Supreme Court affirmed, 4--1, finding no restriction on her religious freedom. The U.S.\ Supreme Court reversed and remanded on June 17, 1963.

\section{Judicial Votes}

\textbf{7--2}, reversed.
\begin{itemize}
\item \textbf{Majority:} Justice Brennan, joined by Chief Justice Warren and Justices Black, Douglas, Clark, and Goldberg.
\item \textbf{Concurring:} Justice Douglas.
\item \textbf{Concurring in the result:} Justice Stewart.
\item \textbf{Dissenting:} Justice Harlan, joined by Justice White.
\end{itemize}

\section{Holding}

A State may not condition unemployment benefits so as to force a worker to choose between a cardinal tenet of her faith and forfeiting benefits, unless a ``compelling state interest'' justifies the substantial infringement and no less restrictive alternative would serve it. Extending benefits does not violate the Establishment Clause. The Court did not reach equal protection, and disclaimed any right to benefits for those whose convictions make them nonproductive.

\section{Analysis of Opinions}

\textbf{Majority (Brennan).} \textit{Burden:} although no criminal penalty applied, an indirect burden that impedes observance is invalid; the choice operated like a fine for Saturday worship. Under \textit{Speiser v.\ Randall}, calling benefits a ``privilege'' does not matter, and the statutory protection for Sunday worshippers compounded the problem. \textit{Scrutiny:} only a compelling interest will do; the State's speculative fear of fraudulent claims was neither raised below nor supported, and it did not show that no alternative could prevent fraud. \textit{Braunfeld v.\ Brown} was distinguishable because a uniform day of rest made exemptions unworkable. \textit{Establishment:} benefits for Sabbatarians reflect ``the governmental obligation of neutrality in the face of religious differences.''

\textbf{Douglas (concurring).} The Clause fences off religious scruples from majoritarian pressure; the harm was interference with conscience, not money.

\textbf{Stewart (concurring in the result).} The Court ``paper[ed] over'' a ``double-barreled dilemma'': its ``positively wooden'' Establishment Clause cases forbid aid to religion, and \textit{Braunfeld}, with a far heavier burden, should be overruled.

\textbf{Harlan, joined by White (dissenting).} The Act relieved only involuntary unemployment, and South Carolina disqualified everyone unavailable for personal reasons, including childcare. The Court effectively overruled \textit{Braunfeld}; a State may accommodate Sabbatarians but is not compelled to.

\section{Examples: Future Applications}

\begin{enumerate}
\item \textbf{Same-Side: Pacifist Worker.} A Jehovah's Witness who quits rather than build tanks may not be denied benefits (\textit{Thomas v.\ Review Board}, 1981), an exception that survived \textit{Smith}.
\item \textbf{Same-Side: Discretionary Exemptions.} A city that grants secular exemptions in its ``sole discretion'' but refuses a Catholic foster agency triggers strict scrutiny (\textit{Fulton}, 2021).
\item \textbf{Opposite-Side: Peyote Ban.} Counselors fired for religious peyote use lose under a neutral, generally applicable criminal law (\textit{Employment Division v.\ Smith}, 1990).
\item \textbf{Opposite-Side: Social Security Taxes.} An Amish employer must pay; a uniform tax system is a compelling interest (\textit{United States v.\ Lee}, 1982).
\item \textbf{Fence-Sitter: Secular Conscience.} A humanist who refuses Saturday shifts to care for a dying parent raises the Stewart--Harlan worry about favoring religious over comparable secular duties.
\end{enumerate}

\section{Critique}

\textbf{Originalist debate.} Philip Hamburger and Justice Scalia argue the Founding did not compel judicial exemptions from general laws and that \textit{Sherbert} invites judges to weigh religious centrality; Michael McConnell and Kurt Lash find historical support for exemptions.

\textbf{Equality concerns.} \textit{Sherbert} helped birth modern strict scrutiny, and Congress restored its test for federal action in RFRA after \textit{Smith}. But Christopher Eisgruber and Lawrence Sager argue that privileging religion over secular commitments undermines equal citizenship and can license exemptions from civil-rights laws.

\section{Key Quotations}

\begin{quote}
``The ruling forces her to choose between following the precepts of her religion and forfeiting benefits, on the one hand, and abandoning one of the precepts of her religion in order to accept work, on the other hand. Governmental imposition of such a choice puts the same kind of burden upon the free exercise of religion as would a fine imposed against appellant for her Saturday worship.''
\end{quote}
\noindent --- Justice Brennan, for the Court, 374 U.S.\ at 404.

\begin{quote}
``It is too late in the day to doubt that the liberties of religion and expression may be infringed by the denial of or placing of conditions upon a benefit or privilege.''
\end{quote}
\noindent --- Justice Brennan, for the Court, 374 U.S.\ at 404.

\begin{quote}
``It is basic that no showing merely of a rational relationship to some colorable state interest would suffice; in this highly sensitive constitutional area, `[o]nly the gravest abuses, endangering paramount interests, give occasion for permissible limitation,' \textit{Thomas v.\ Collins}, 323 U.S.\ 516, 530.''
\end{quote}
\noindent --- Justice Brennan, for the Court, 374 U.S.\ at 406.

\end{document}
